% from main.tex
%%%%%%%%%%  scan idx 170  |  folio 157  %%%%%%%%%%
large,row sep=large]
E^\vee\otL F\otL F^\vee\otL E \arrow[r,"\Id\otimes\tau_E"] \arrow[d,"\Id\otimes\tau_F"'] & F^\vee\otL F \arrow[d,"\tau_F"]\\
E^\vee\otL E \arrow[r,"\tau_E"'] & A .
\end{tikzcd}
\]

\textbf{Corollaire 8.4.} \(E,F,U\) étant donnés comme en (8.2), soient \(u\in\Hom(E|U,E|U)\), \(s:E|U\to F|U\) un isomorphisme, d'où \(sus^{-1}\in\Hom(F|U,F|U)\). On a alors
\[
  \Tr_{F|U}(sus^{-1})=\Tr_{E|U}(u).
\]

\emph{Preuve.} C'est immédiat.

\textbf{Remarque 8.5.} Si \(\mathcal C\) est une \(S\)-catégorie fibrée, notons \(\mathcal E(\mathcal C)\) la \(S\)-catégorie fibrée dont la fibre en \(U\in\Ob S\) est la catégorie ayant pour objets les couples \((L,u)\) où \(L\in\Ob\mathcal C_U\) et \(u\in\operatorname{End}_S(L)\) et pour flèches de \((L,u)\) à \((M,v)\) les flèches \(f:L\to M\) telles que \(fu=vf\). On peut, d'après (8.4), considérer l'homomorphisme trace comme un \(S\)-foncteur de \(\mathcal E(\Dpx(S)_{\parf})^{\mathrm{is}}\) dans \(\mathcal A\) (où l'exposant \(\mathrm{is}\) signifie que les seules flèches considérées sont les isomorphismes, et \(\mathcal A\) est la \(S\)-catégorie à fibres discrètes définie par \(A\) (cf. (6.6.a))). Comme tout complexe parfait est localement de tor-dimension finie, le \(S\)-foncteur d'oubli
\[
  \mathcal E(\Dpx(S)_{\parf})^{\mathrm{is}}
  \longrightarrow
  \mathcal E(\Dpx(S)_{\parf})^{\mathrm{is}}
\]
induit un isomorphisme sur les champs associés. Par suite, l'homomorphisme
